% !TEX program = lualatex
\documentclass[11pt,a4paper]{article}
\usepackage{icat-common}

\theoremstyle{definition}
\newtheorem{definition}{定義}[section]
\newtheorem{axiom}{公理}[section]
\newtheorem{lemma}{補題}[section]
\newtheorem{theorem}{定理}[section]
\newtheorem{corollary}{系}[section]
\newtheorem{hypothesis}{仮説}[section]
\newtheorem{remark}{備考}[section]

\newcommand{\Einfty}{E_\infty}
\newcommand{\SelfRef}{\Sigma}
\renewcommand{\Pr}{\mathrm{Pr}}

\title{観測者の極限消去：\\
ゲーデルの決定不能性を不動点として回収する操作について}
\author{千葉将臣}
\date{2026-06-19}

\begin{document}
\maketitle

\begin{abstract}
ゲーデルの第一不完全性定理が生成する決定不能文 $G_T$ は、
固定された体系（観測者）$T$ に依存して現れる。
本稿は、極限操作が観測者を消去する装置である理由を
顕現論の消去公理から原理的に導出し、
極限によって観測者を消去した後に残る不動点 $\Einfty$ を定式化する。
$G_T$ は $\Einfty$ の観測者依存的な断面に過ぎないという仮説を提示し、
決定不能性とは観測者の固定が生んだ見かけの壁であることを論証する。
\end{abstract}

%=============================================================
\section{観測者としての形式体系}
%=============================================================

ゲーデルの第一不完全性定理は以下を主張する。
一貫した再帰的公理化可能な体系 $T$（ロビンソン算術 $Q$ を内包）に対し、
$$T \nvdash G_T \quad \text{かつ} \quad T \nvdash \neg G_T$$
を満たす文 $G_T$ が存在する。

ここで $T$ は\textbf{観測者}として機能している。$G_T$ は $T$ という
固定された視点から見たときに初めて「決定不能」として現れる。
$T$ を別の体系 $T'$（$T' \supset T$）に移せば $G_T$ は証明可能になるが、
今度は $G_{T'}$ が新たに現れる。決定不能性は観測者の移動に追随する。

\begin{definition}[観測者依存性]
文 $\varphi$ が\textbf{観測者依存的}であるとは、ある体系 $T$ に対して
$T \nvdash \varphi$ かつ $T' \vdash \varphi$（$T' \supsetneq T$）
を満たす体系 $T'$ が存在することをいう。
$G_T$ は観測者依存的である。
\end{definition}

%=============================================================
\section{消去公理と極限操作の原理的接続}
%=============================================================

極限操作が観測者を消去する装置である理由は、
単なる記法上の偶然ではない。
顕現論の消去公理がその原理的根拠を与える。

\begin{axiom}[消去公理 \cite{Chiba2026_chikaku} ]
消去の収束は確定した値として完了する。
この完了した値としてある概念の極限・無限が生ずる。
\end{axiom}

消去公理の前段として、顕現論は以下を定義している。

\begin{definition}[消去 \cite{Chiba2026_chikaku} ]
「ある」の知覚と「ない」の認識の同時成立により、
「ない」の持つ高次元の可能性が低次元の「ある」として
捕捉可能な形に収束する操作を\textbf{消去}と呼ぶ。
\end{definition}

\begin{definition}[世界の立ち上がり \cite{Chiba2026_chikaku} ]
世界が「ある」として知覚可能な状態へ移行する際、
「ない」の高次元性と極限の確定が同時に生じる。
この同時性が消去の発生を可能にする機序である。
\end{definition}

消去公理から極限操作の観測者不在性を導出する。

\begin{lemma}[極限操作の観測者不在性]
消去公理（公理2.1）より、極限操作は観測者を消去する装置である。
すなわち、極限 $\lim_{n \to \infty} a_n = L$ のε-δ定義：
$$\forall \varepsilon > 0,\ \exists N \in \mathbb{N},\ \forall n > N,\
|a_n - L| < \varepsilon$$
には特定の観測者への参照が存在せず、
極限値 $L$ は観測者から独立した対象として確定する。
\end{lemma}

\begin{proof}
消去公理は「消去の収束が確定した値として完了する」ことを保証する。
ε-δ定義は全称文と存在文の連鎖であり、
「誰が極限を取るか」を構造的に要求しない。
これは消去公理が記述する「完了した値」の形式的実現である——
観測者（「ある」を捕捉する主体）を消去した後に、
収束の完了した値 $L$ が確定する。
したがって極限操作は消去公理の数学的実装として、
観測者不在の装置として機能する。
\end{proof}

\begin{remark}
数学の実践において極限は「誰かが取る操作」として運用されてきたが、
定義そのものに観測者は不在である。
この自明だが見落とされてきた事実は、
消去公理を参照することで「なぜ不在なのか」の原理的根拠を持つ。
\end{remark}

%=============================================================
\section{観測者消去後の不動点 $\Einfty$}
%=============================================================

対角補題を $\varphi(x) := \neg \Pr_T(x)$ に適用すると、
$$T \vdash G \leftrightarrow \neg \Pr_T(\ulcorner G \urcorner)$$
を満たす文 $G$ が $T$ の内部で構成される。
自己言及演算子 $\SelfRef_T(A) := \neg \Pr_T(\ulcorner A \urcorner)$ を定義すると、
$G$ は $\SelfRef_T$ の不動点方程式：
$$G = \SelfRef_T(G)$$
を満たす。この等式の $=$ は $\lim_{n\to\infty} a_n = L$ と同型の
静的記述であり、プロセスの収束先を観測者なしに確定させる。
これは補題2.1が保証する極限操作の観測者不在性の直接の帰結である。

\begin{definition}[極限同値 $\Einfty$]
$\SelfRef_T$ の不動点を\textbf{極限同値}と呼び $\Einfty$ と表記する。
$$\Einfty = \SelfRef_T(\Einfty)
\quad \Longleftrightarrow \quad
T \vdash \Einfty \leftrightarrow \neg \Pr_T(\ulcorner \Einfty \urcorner)$$
\end{definition}

\begin{theorem}[観測者消去による $G_T$ の回収]
観測者（体系 $T$）を固定した場合、$\Einfty$ は決定不能文 $G_T$ として現れる。
観測者を極限操作によって消去した場合、同一の対象が不動点 $\Einfty$ として確定する。
\end{theorem}

\begin{proof}
対角補題により $\Einfty$ は $T$ の内部で構成される。
$T$ を固定した観測者として採用した場合、$T$ の一貫性のもとで
$T \nvdash \Einfty$ かつ $T \nvdash \neg \Einfty$ が成立する（ゲーデルの定理）。
これが $G_T$ である。

一方、$\Einfty = \SelfRef_T(\Einfty)$ は $T$ への参照を含むが、
この等式の成立自体は $T$ を観測者として固定することを要求しない——
不動点方程式はその解が存在することを観測者なしに記述する（補題2.1）。
したがって $\Einfty$ は観測者消去後も確定した対象として残る。
\end{proof}

定理3.1は「観測者 $T$ を固定した場合 $G_T$ として現れ、
消去後は $\Einfty$ として確定する」という構造を示す。
次の仮説はその関係をより強く述べるものであり、
現時点では証明を持たない。

\begin{hypothesis}[決定不能文の共通断面]
任意の観測者 $T$ に対して生じる決定不能文 $G_T$ は、
共通対象 $\Einfty$ の観測者依存的断面である：
$$G_T \mweq \Einfty \text{ の } T\text{-断面}$$
すなわち観測者の多様性は $\Einfty$ という単一対象への
異なる切り口として統一的に記述できる。
\end{hypothesis}

\begin{remark}
$G_T$ が $\Einfty$ の観測者依存的断面であることは、
$\Einfty$ を観測者消去後の確定した値として
天下りに与えることの帰結として読まれるべきであり、
体系内での証明を持たない。
これは真理保存性が体系内で証明できないことと同型の構造である。
\end{remark}

\begin{remark}[仮説の動機]
各観測者 $T$ に対して決定不能文が繰り返し現れるという現象は、
背後に同一の構造があることを示唆する。
これは観測者ごとに独立した現象ではなく、
$\Einfty$ という共通対象の異なる断面として
統一的に読める可能性を持つ。
この発想は数学的証明というより物理学的直観——
繰り返し現れる現象は共通の構造を持つ——に近い。
仮説4.1の証明は今後の課題である。
\end{remark}

%=============================================================
\section{帰結：数学の有用性の起源}
%=============================================================

数学がウィグナーの言う「不合理なほどの有効性」を物理世界に対して持つ理由は、
数学的構造が観測者消去後に残る構造だからである。
極限・積分・微分方程式はいずれも消去公理の帰結として
観測者を消去した記述であり、
物理法則もまた観測者（「今」「ここ」）を含まない。

$\Einfty$ はこの構造の論理的核心に位置する。
自己言及という観測者に最も依存しそうな操作が、
消去公理に従った極限によって観測者から解放されたとき、
不動点という最も安定した数学的対象に変わる。

決定不能性は限界ではない。観測者を消し忘れた記述の痕跡である。

\begin{thebibliography}{9}

\bibitem{Chiba2026_chikaku}
千葉将臣.
\newblock 顕現論における知覚原理.
\newblock \emph{空数学・界論基礎論考}, 2026.

\bibitem{Chiba2026_einfty}
千葉将臣.
\newblock 真理保存性を持つ体系が必然的に $E_\infty$ を含む：
対角補題による極限同値の内在的定式化.
\newblock \emph{空数学・界論基礎論考}, 2026.

\end{thebibliography}

\end{document}
